\section{Cơ sở chung}
\subsection{Đồ thị tri thức và đồ thị tri thức đa phương thức}

Đồ thị tri thức (Knowledge Graph - KG) biểu diễn tri thức dưới dạng đồ thị,
trong đó các thực thể (entity) là đỉnh và các sự kiện là các cạnh được gắn nhãn quan hệ (relation) kết nối các thực thể với nhau \cite{balkir2021}. 

Thực thể (entity) chính là một đối tượng nào đó. Chẳng hạn như ``Trường Đại học Khoa học Tự nhiên", ``Khoa Công nghệ thông tin", ``Môn học Xử lý ngôn ngữ tự nhiên",...
 Mỗi thực thể được biểu diễn bằng một node trên đồ thị. Quan hệ (relation) là sự mô tả liên kết giữa các thực thể.
 Chẳng hạn như ``được giảng dạy bởi'', ``là đơn vị của'',... Quan hệ được biểu diễn bằng cạnh có hướng và có nhãn trên đồ thị.
 Ví dụ hai node ``Môn học Xử lý ngôn ngữ tự nhiên'' và ``Khoa Công nghệ thông tin'' sẽ được nối bằng quan hệ ``được giảng dạy bởi''.
 Hướng của cạnh cho biết rằng ``Môn học Xử lý ngôn ngữ tự nhiên \textit{được giảng dạy bởi} Khoa Công nghệ thông tin''.

Đồ thị tri thức được ký hiệu $\mathcal{G} = (\mathcal{E}, \mathcal{R}, \mathcal{F})$,
trong đó $\mathcal{E}$ là tập các đỉnh đại diện cho các thực thể đầu $h$ (head) và đuôi $t$ (tail),
$\mathcal{R}$ là tập các loại quan hệ $r$ được biểu diễn bằng nhãn của các cạnh, còn $\mathcal{F} \subseteq \mathcal{E} \times \mathcal{R} \times \mathcal{E}$ là tập các bộ ba $(h,r,t)$.


Từ ví dụ trên, thông tin về ``Khoa Công nghệ thông tin'' còn có thể được biểu diễn qua nhiều loại dữ liệu như: cấu trúc đồ thị (``là đơn vị của Trường Đại học Khoa học Tự nhiên''),
\textit{văn bản} (miêu tả về khoa), \textit{hình ảnh} (hình ảnh được gắn với khoa), \textit{thuộc tính} (số ngành đào tạo, số giảng viên,...).
Từ đây ta dẫn đến khái niệm đồ thị tri thức đa phương thức (Multimodal Knowledge Graph - MMKG).
Cụ thể, đồ thị tri thức đa phương thức kết hợp cấu trúc quan hệ với các nguồn dữ liệu khác như hình ảnh, văn bản,... để cung cấp thêm thông tin về thực thể.
Mỗi nguồn thông tin thể hiện một khía cạnh của thực thể và không nhất thiết thực thể nào cũng có đủ các nguồn thông tin này.

Để mô hình có thể xử lý được, mỗi thực thể được biểu diễn bằng một vector số thực.
Vector này có thể được xây dựng từ nhiều nguồn thông tin và kết hợp thành biểu diễn tổng hợp.
Cách mã hóa và kết hợp sẽ được trình bày trong phần phương pháp của từng bài báo.


\subsection{Hai bài toán MMEA và MKGR}

Hai bài báo trong báo cáo này giải quyết hai bài toán hoàn toàn khác nhau.

Multimodal Entity Alignment - MMEA (căn chỉnh thực thể đa phương thức) là bài toán mà công trình CateEA giải quyết. 
rong thực tế, có nhiều đồ thị tri thức được xây dựng độc lập từ các nguồn khác nhau.
Chẳng hạn một đồ thị tiếng Việt và một đồ thị tiếng Anh, hai đồ thị lần lượt có thực thể 
\textit{``Khoa Công nghệ thông tin - Trường Đại học Khoa học Tự nhiên''} và \textit{``Faculty of Information Technology - HCMUS''} cùng chỉ một khoa của một trường,
nhưng được biểu diễn bằng hai tên khác nhau và không có liên kết trực tiếp.
Bài toán MMEA đặt ra câu hỏi: \textit{Cho trước hai đồ thị tri thức đa phương thức và một số cặp thực thể tương ứng đã biết, làm sao để tìm ra các cặp thực thể tương ứng còn lại?}
Đầu ra của bài toán là danh sách các cặp thực thể từ hai đồ thị cùng chỉ về một đối tượng ngoài đời thực.

Multimodal Knowledge Graph Reasoning - MKGR (suy luận trên đồ thị tri thức đa phương thức) là bài toán mà công trình DSoM giải quyết.
Trên thực tế, hầu hết các đồ thị tri thức đều không hoàn chỉnh, sẽ có những thông tin bị thiếu trong đồ thị.
Bài toán MKGR đặt ra câu hỏi: Cho trước một bộ ba bị khuyết một thành phần, hãy dự đoán thành phần đó.
Ví dụ với bộ ba \textit{(Môn học Xử lý ngôn ngữ tự nhiên, được giảng dạy bởi, ?)}, mô hình cần suy luận để điền vào thông tin thiếu là \textit{``Khoa Công nghệ thông tin''}.
Đầu ra là danh sách các thực thể ứng viên được sắp xếp theo độ tin cậy. Trong công trình DSoM, tác giả chỉ xét dự đoán $(?,r,t)$ và $(h,r,?)$, không dự đoán quan hệ $(h,?,t)$.